\SourceLeaf{422}%
\begin{center}6.\end{center}
\noindent\hspace*{1.2em}Wenn man $x$ und $y$ als Functionen von $s$ betrachtet, so fallen die Formeln noch zierlicher aus. Man setze
\[
\begin{aligned}
\frac{dx}{ds}&=x',&\qquad \frac{dx'}{ds}&=x'',\\
\frac{dy}{ds}&=y',&\frac{dy'}{ds}&=y''.
\end{aligned}
\]
Es ist dann
\[
\begin{aligned}
x'&=\cos\varphi,&\qquad y'&=\sin\varphi,\\
x''&=-\frac{\sin\varphi}{r},&y''&=\frac{\cos\varphi}{r},
\end{aligned}
\]
oder
\[
y'=-rx'',\qquad x'=ry''
\]
oder auch
\[
1=r(x'y''-y'x''),
\]
so dass
\[
x'y''-y'x''
\]
die Krümmung und
\[
\frac{1}{x'y''-y'x''}
\]
den Krümmungshalbmesser vorstellt.

\begin{center}7.\end{center}
\noindent\hspace*{1.2em}Wir schreiten jetzt zur Betrachtung der krummen Flächen fort. Zur Versinnlichung der Richtungen von geraden Linien im Raume nach seinen drei Dimensionen betrachtet denken wir uns um einen beliebigen Mittelpunkt eine Kugelfläche mit dem Halbmesser $1$ beschrieben; ein Punkt derselben repräsentirt demnach die Richtung aller mit dem Radius, an dessen Endpunkt er liegt, parallelen geraden Linien. So wie die Lagen aller Punkte im Raume durch die senkrechten Abstände $x,y,z$ von dreien auf einander senkrechten Ebenen bestimmt werden, so soll die Richtung der drei Hauptaxen, welche auf diese Hauptebenen senkrecht sind, auf der Oberfläche der Hülfskugel
